\ncdlesson{5}{Hvornår går bussen?}{公交车几点开？}
{DU3 Modul 1（A1+）｜对标 NCE 1 Lesson 21--26}
{钟点表达（\textit{klokken..., halv..., kvart i/over..., minutter i/over...}）与星期表达（\textit{på mandag} vs \textit{om mandagen}）· 时间疑问词（\textit{hvornår, hvor tit, hvor lang tid}）· 句首时间/地点状语引发的主句 V2 倒装}
{本课进入第二单元“作息与交通”，同时迎来丹麦语最重要的一条句法铁律——\textbf{V2 倒装}！请重点攻克“\textit{halv ni} 是八点半而不是九点半”的思维切换，并把 \textit{Om morgenen kører bussen...} 练成脱口而出的条件反射。}

% ---------------------------------------------------------------------
\begin{foerduleaeser}
带着下面两个问题阅读课文对话：
\begin{enumerate}
  \item \textit{Hvor tit kører bus 4A om morgenen, og hvor tit kører den i weekenden?}（4A 路公交车早上多久发一班？周末又多久一班？）
  \item \textit{Hvornår møder Emma på skolen om mandagen, og hvor lang tid tager turen til banegården?}（艾玛每周一几点到校？去中心火车站的车程要多久？）
\end{enumerate}
\end{foerduleaeser}

% ---------------------------------------------------------------------
\ncdsubsec{二、Dialog \& Udtale（课文对话、发音标注与中文译文）}

\begin{tipbox}{场景说明}
周一早上八点零五分，Emma 站在公交站台查看电子站牌，向一旁同样等车的本地通勤者 Henrik 打听 4A 路公交车的发车时间与车程。
\end{tipbox}

\dialogturn{1}{Emma}{Undskyld, hvornår går bus 4A mod centrum?}
{\textit{hvornår} \ipa{vɒˈnɒːˀ}（词首 \textbf{h 不发音}）；\textit{går} \ipa{ˈɡɒːˀ}（长圆唇音 + 喉塞音）；\textit{mod} \ipa{ˈmoːðˀ}（软 \ipa{ð} + 喉塞音）。}
{打扰一下，开往市中心的 4A 路公交车几点开？}

\dialogturn{2}{Henrik}{Lige nu er klokken fem minutter over otte. Den næste bus går kvart over otte.}
{\textit{klokken} \ipa{ˈkʰlʌɡən}（口语常略去中间元音读如 \ipa{kʰlʌɡŋ}）；\textit{fem minutter over} \ipa{ˈfɛmˀ miˈnudɐ ˈʌʊ̯ˀɐ}；\textit{kvart} \ipa{ˈkvɑːd}（\textit{r} 融进长元音 \ipa{ɑː}）。}
{现在是八点零五分。下一班车八点一刻开。}

\dialogturn{3}{Emma}{Super! Hvor tit kører bussen her om morgenen?}
{\textit{hvor tit} \ipa{vɒ ˈtʰid}（频率疑问“多久一次”）；\textit{kører} \ipa{ˈkʰøːɐ}（\textit{r} 弱化融尾）；\textit{om morgenen} \ipa{ʌm ˈmɒːnən}（\textbf{g 完全不发音}！）。}
{太好了！这儿的公交车早上多久发一班？}

\dialogturn{4}{Henrik}{Om morgenen kører den hvert tiende minut, men i weekenden kører den kun én gang i timen.}
{\textit{tiende} \ipa{ˈtʰiːənə}（\textbf{d 不发音}）；\textit{hvert} \ipa{ˈvɛɐ̯ˀd}（\textbf{h 不发音}）；\textit{weekenden} \ipa{ˈviːɡɛndən}；\textit{én gang} \ipa{ˈeːˀn ˈɡaŋˀ}（\textit{én} 带重音表示数量“一”）。}
{早上每十分钟一班，但周末每小时才一班。}

\dialogturn{5}{Emma}{Okay, det er godt at vide. Hvor lang tid tager turen ind til banegården?}
{\textit{at vide} \ipa{ad ˈviːðə}（软 \ipa{ð}）；\textit{hvor lang tid} \ipa{vɒ ˈlaŋˀ ˈtʰiːðˀ}；\textit{banegården} \ipa{ˈbæːnəˌɡɒːˀən}（\textit{gård} 的 \textbf{d 不发音}）。}
{好的，知道这个太有用了。坐到中心火车站需要多长时间？}

\dialogturn{6}{Henrik}{Herfra tager det cirka tyve minutter. Skal du på sprogskole eller på arbejde?}
{\textit{herfra} \ipa{ˈheɐ̯ˀˌfʁɑ}；\textit{cirka} \ipa{ˈsiɐ̯ɡa}；\textit{arbejde} \ipa{ˈɑːbɑɪ̯ˀdə}（\textit{ar-} 读长音 \ipa{ɑː}，\textit{-ej-} 读 \ipa{ɑɪ̯}）。}
{从这儿出发大约要二十分钟。你是要去语言学校还是去上班？}

\dialogturn{7}{Emma}{Jeg skal på arbejde i dag. Men på mandag starter mit nye hold på skolen.}
{\textit{på arbejde} \ipa{pʰɔ ˈɑːbɑɪ̯ˀdə}；\textit{mandag} \ipa{ˈmanˀda}（\textbf{nd 的 d 和 dag 的 g 均不发音}！）。}
{我今天去上班。不过下周一我的新班就要在学校开课了。}

\dialogturn{8}{Henrik}{Spændende! Hvornår møder du på skolen om mandagen?}
{\textit{spændende} \ipa{ˈsbɛnənə}（\textbf{nd 中的 d 不发音}，常简读为 \ipa{ˈsbɛnə}）；\textit{om mandagen} \ipa{ʌm ˈmanˀdæːˀən}（加 \textit{-en} 表示“每周一”）。}
{真让人期待！你每周一几点到学校上课？}

\dialogturn{9}{Emma}{Om mandagen møder jeg klokken halv ni, så jeg tager bussen fem minutter i otte. Men på fredag har jeg et tidligt møde kvart i otte!}
{\textit{halv ni} \ipa{ˈhalˀ ˈniːˀ}（重要！\textbf{v 永远不发音}，\textit{l} 带喉塞音）；\textit{fredag} \ipa{ˈfʁɛːˀda}（\textbf{g 静音}）；\textit{tidligt} \ipa{ˈtʰiːðli}（\textbf{g 和 t 均静音}）。}
{每周一我八点半上课，所以我七点五十五分坐公交车。不过这周五我七点四十五分有个早会！}

\dialogturn{10}{Henrik}{Se, nu kommer bussen faktisk! Klokken er ti minutter over otte præcis. God arbejdsdag!}
{\textit{faktisk} \ipa{ˈfaɡdisg}；\textit{præcis} \ipa{pʰʁɛˈsiːˀs}；\textit{arbejdsdag} \ipa{ˈɑːbɑɪ̯dsˌdæːˀ}（\textit{d} 音极弱，整体听起来像 \ipa{ˈɑːbɑɪ̯sˌdæːˀ}）。}
{看，公交车这不就来了！正好八点十分整。祝你工作顺利！}

% ---------------------------------------------------------------------
\ncdsubsec{三、Gloser（核心生词表）}

\begin{tipbox}{词汇提示}
位移动词 \textit{at gå} 和 \textit{at køre} 在现在完成时中，若表示\textbf{明确移向某地/到达终点}，助动词用 \dabold{er}；若仅强调持续动作本身（如兜风、开了很久车），才用 \textit{har}。
\end{tipbox}

\begin{glosertable}
疑问词 & hvornår & --- & 什么时候（问时间点） \\
疑问词 & hvor tit / hvor lang tid & --- & 多久一次（频率）/ 多长时间（时长） \\
\midrule
名词 & en bus & \textit{(bussen, busser, busserne)} & 公交车（单数定指与复数双写 \textit{s}） \\
名词 & en klokke & \textit{(klokken, klokker, klokkerne)} & 钟表；\textit{klokken...} ……点钟 \\
名词 & et kvarter & \textit{(kvarteret, kvarter / kvarterer, kvartererne)} & 一刻钟（15 分钟） \\
名词 & en tur & \textit{(turen, ture, turene)} & 一趟行程；散步，出行 \\
名词 & en banegård & \textit{(banegården, banegårde, banegårdene)} & 火车站（主站） \\
名词 & et stoppested & \textit{(stoppestedet, stoppesteder, stoppestederne)} & 公交站台（\textit{d} 常弱化） \\
名词 & et arbejde & \textit{(arbejdet, arbejder, arbejderne)} & 工作（\textit{på arbejde} 上班） \\
名词 & en weekend & \textit{(weekenden, weekender, weekenderne)} & 周末（\textit{i weekenden} 在周末） \\
名词 & en morgen & \textit{(morgenen, morgener, morgenerne)} & 早晨（\textit{om morgenen} 每天早上） \\
名词 & et møde & \textit{(mødet, møder, møderne)} & 会议；会面 \\
\midrule
动词 & at vide & \textit{(ved, vidste, har vidst)} & 知道，了解（高频不规则变化） \\
动词 \dabold{[er]} & at gå & \textit{(går, gik, \dabold{er gået} / har gået)} & 走；（车船）发车开行（到达终点用 \textit{er}） \\
动词 \dabold{[er]} & at køre & \textit{(kører, kørte, \dabold{er kørt} / har kørt)} & 驾驶；行驶（去往某地用 \textit{er}） \\
动词 & at starte & \textit{(starter, startede, har startet)} & 开始；启动 \\
动词 & at møde & \textit{(møder, mødte, har mødt)} & 到岗，报道上班/上课；遇见 \\
\midrule
副/形容词 & cirka / næste / tidlig & \textit{cirka (不变)} / \textit{næste (不变)} / \textit{tidlig (-t, -e)} & 大约 / 下一个的 / 早的 \\
\end{glosertable}

% ---------------------------------------------------------------------
\ncdsubsec{四、Grammatik \& Faste vendinger（核心语法与地道表达精讲）}

\begin{grammatikbox}{语法点 1：丹麦语钟点表达（Klokken）——警惕 \textit{halv} 陷阱！}
询问时间：\da{Hvad er klokken?}（现在几点了？）/ 回答：\da{Klokken er...}。其中“半点”看的是\textbf{下一个整点}，这是中国学生最常看错的考点：

\vspace{0.25em}
\renewcommand{\arraystretch}{1.24}
{\small
\begin{tabularx}{\linewidth}{@{} >{\bfseries\color{NCDRed}}p{1.6cm} >{\itshape\color{NCDNavy}\raggedright\arraybackslash}p{5.8cm} >{\raggedright\arraybackslash}X @{}}
\toprule
\textbf{时间} & \textbf{丹麦语地道说法} & \textbf{记忆逻辑} \\
\midrule
08:00 & klokken otte & 整点直接报数（常用 12 小时制） \\
08:10 & ti minutter \dabold{over} otte & 过了 8 点 10 分（前半小时用 \textit{over}） \\
08:15 & \dabold{kvart over} otte & 过 8 点一刻 \\
08:25 & fem minutter \dabold{i halv ni} & 差 5 分到“八点半（半九）” \\
08:30 & \dabold{halv ni} & \textbf{不是 9:30！}是“走向 9 点的一半”= 8:30 \\
08:35 & fem minutter \dabold{over halv ni} & 过“八点半”5 分钟 \\
08:45 & \dabold{kvart i} ni & 差一刻到 9 点（后半小时用 \textit{i}） \\
08:55 & fem minutter \dabold{i} ni & 差 5 分到 9 点 \\
\bottomrule
\end{tabularx}
}
\end{grammatikbox}

\begin{grammatikbox}{语法点 2：星期表达（Ugedage）与时间介词 \textit{på} vs \textit{om}}
丹麦语星期名称\textbf{首字母全部小写}：\textit{mandag, tirsdag, onsdag, torsdag, fredag, lørdag, søndag}。前面用什么介词，完全取决于“一次性”还是“规律性重复”：
\begin{itemize}
  \item \dabold{på + 星期原形} = \textbf{即将到来的那一个星期几}（单次、将来）：\newline \textit{\da{På mandag} starter mit nye hold.}（下周一我的新班开课。）/ \textit{\da{På fredag} har jeg et møde.}（这周五我有个会。）
  \item \dabold{om + 星期 + -en} = \textbf{每周几（习惯性规律重复）}：\newline \textit{\da{Om mandagen} møder jeg klokken halv ni.}（我每周一都八点半上课。）同理：\da{om morgenen}（每天早上），\da{om aftenen}（每天晚上），\da{om dagen}（白天）。
  \item \textbf{特殊搭配}：\textit{\da{i dag}}（今天），\textit{\da{i weekenden}}（在周末），\textit{\da{i timen}}（每小时）。
\end{itemize}
\end{grammatikbox}

\begin{grammatikbox}{语法点 3：时间疑问词三剑客（\textit{hvornår / hvor tit / hvor lang tid}）}
\begin{itemize}
  \item \da{Hvornår...?}（什么时候/几点？）$\rightarrow$ 问具体时间点。\newline 例：\textit{Hvornår går bussen?} $\rightarrow$ \textit{Den går kvart over otte.}
  \item \da{Hvor tit...?}（多久一次？）$\rightarrow$ 问发生频率。\newline 例：\textit{Hvor tit kører den?} $\rightarrow$ \textit{Hvert tiende minut / én gang i timen / to gange om ugen.}
  \item \da{Hvor lang tid...?}（多长时间？）$\rightarrow$ 问持续时长。\newline 例：\textit{Hvor lang tid tager turen?} $\rightarrow$ \textit{Det tager cirka tyve minutter.}
\end{itemize}
\end{grammatikbox}

\begin{grammatikbox}{语法点 4：丹麦语核心铁律——句首时间/地点状语触发主句 V2 倒装！}
V2（Verb Second）是指\textbf{任何丹麦语陈述主句中，变位动词必须永远占据第 2 个句子成分的位置}！只要句首不再是主语，而是放了时间或地点状语，\textbf{主语就必须立刻退到动词后面}：

\vspace{0.25em}
\renewcommand{\arraystretch}{1.24}
{\small
\begin{tabularx}{\linewidth}{@{} >{\bfseries\color{NCDGray}}p{3.0cm} >{\raggedright\arraybackslash}p{4.0cm} >{\bfseries\color{NCDRed}}p{2.3cm} >{\raggedright\arraybackslash}X @{}}
\toprule
\textbf{句型} & \textbf{第 1 位（主题/状语）} & \textbf{第 2 位（动词）} & \textbf{第 3 位（主语）+ 其他} \\
\midrule
正常语序（SVO） & \da{Den}（主语） & kører & hvert tiende minut om morgenen. \\
时间状语在句首 & \textbf{Om morgenen} & kører & \da{den} hvert tiende minut. \\
时间状语在句首 & \textbf{På mandag} & starter & \da{mit nye hold} på skolen. \\
地点状语在句首 & \textbf{Herfra} & tager & \da{det} cirka tyve minutter. \\
\bottomrule
\end{tabularx}
}
\vspace{0.25em}

\textbf{避坑警告}：千万别按英语语序说成 *\textit{Om morgenen den kører...}！一旦逗号或时间词后面直接跟了 \textit{den / jeg / bussen}，语法就错了。
\end{grammatikbox}

% ---------------------------------------------------------------------
\ncdsubsec{五、Mønsterøvelser（新概念式句型操练）}

\textbf{练习 A：用丹麦语写出下列钟点（不用阿拉伯数字）}\\[0.3em]
\begin{tabularx}{\linewidth}{@{} X X @{}}
1. 07:00 $\rightarrow$ \underline{\hspace{4.2cm}} & 4. 08:30 $\rightarrow$ \underline{\hspace{4.2cm}} \\
2. 08:15 $\rightarrow$ \underline{\hspace{4.2cm}} & 5. 08:45 $\rightarrow$ \underline{\hspace{4.2cm}} \\
3. 08:25 $\rightarrow$ \underline{\hspace{4.2cm}} & 6. 09:50 $\rightarrow$ \underline{\hspace{4.2cm}} \\
\end{tabularx}
\vspace{0.3em}

\textbf{练习 B：将句尾的时间或地点状语提前到句首，按 V2 倒装规则改写句子}
\begin{quote}
\small 例：\textit{Bussen kører hvert tiende minut \textbf{om morgenen}.} $\rightarrow$ \textit{\textbf{Om morgenen} kører bussen hvert tiende minut.}
\end{quote}
\begin{enumerate}
  \item \textit{Turen tager cirka tyve minutter \textbf{herfra}.} $\rightarrow$
  \item \textit{Jeg starter på et nyt hold \textbf{på mandag}.} $\rightarrow$
  \item \textit{Bussen kører kun én gang i timen \textbf{i weekenden}.} $\rightarrow$
  \item \textit{Vi møder på arbejde klokken ni \textbf{om tirsdagen}.} $\rightarrow$
  \item \textit{Emma venter på bus 4A \textbf{ved stoppestedet}.} $\rightarrow$
\end{enumerate}

\textbf{练习 C：选填 \textit{på / om / hvor tit / hvor lang tid / hvornår} 及中译丹}
\begin{enumerate}
  \item \underline{\hspace{2.8cm}} (\textit{Hvornår / Hvor lang tid}) tager turen med tog til Aarhus?
  \item \underline{\hspace{2.8cm}} (\textit{Hvor tit / Hvornår}) går du på sprogskole? --- To gange om ugen.
  \item Jeg har møde \underline{\hspace{1.2cm}} (\textit{på / om}) onsdag, men jeg arbejder hjemme \underline{\hspace{1.2cm}} (\textit{på / om}) fredagen.
  \item 每天早上公交车七点半从火车站发车。（用 \textit{Om morgenen} 开头）
  \item 现在几点了？--- 现在差五分八点半（8:25）。
\end{enumerate}

% ---------------------------------------------------------------------
\ncdsubsec{六、Resumé \& Mundtlig træning（口语复述与参考答案）}

\textbf{口语复述任务}：遮住课文，先口头回答 \textbf{Før du læser} 的两个问题，然后以“\textit{Om mandagen... / På fredag...}”做至少 3 句倒装，向同伴介绍你常用的公交车或地铁发车频率与上学/上班时间。

\begin{answerbox}
\textbf{练习 A 答案}\\[0.2em]
\begin{tabularx}{\linewidth}{@{} X X @{}}
1. \da{klokken syv} & 4. \da{halv ni}（注意不是 \textit{halv otte}！） \\
2. \da{kvart over otte} & 5. \da{kvart i ni} \\
3. \da{fem minutter i halv ni} & 6. \da{ti minutter i ti} \\
\end{tabularx}
\vspace{0.35em}

\textbf{练习 B 答案}
\begin{enumerate}
  \item \textit{Herfra \dabold{tager det} cirka tyve minutter.}
  \item \textit{På mandag \dabold{starter jeg} på et nyt hold.}
  \item \textit{I weekenden \dabold{kører bussen} kun én gang i timen.}
  \item \textit{Om tirsdagen \dabold{møder vi} på arbejde klokken ni.}
  \item \textit{Ved stoppestedet \dabold{venter Emma} på bus 4A.}
\end{enumerate}

\textbf{练习 C 答案}
\begin{enumerate}
  \item \da{Hvor lang tid}（问车程耗时）
  \item \da{Hvor tit}（回答是“每周两次”，问频率）
  \item \da{på}（下周三/这周三，单次）；\da{om}（每周五，规律）
  \item \textit{Om morgenen \dabold{går bussen} fra banegården klokken halv otte.}
  \item \textit{Hvad er klokken nu? --- Klokken er fem minutter i halv ni.}
\end{enumerate}

\textbf{Før du læser 参考回答 \& 口语复述范文（Resumé）}
\begin{itemize}
  \item \textbf{问题 1}：\textit{Om morgenen kører bus 4A hvert tiende minut, men i weekenden kører den kun én gang i timen.}
  \item \textbf{问题 2}：\textit{Om mandagen møder Emma på skolen klokken halv ni. Turen ind til banegården tager cirka tyve minutter.}
  \item \textbf{口语复述范文}：\textit{Emma venter på bus 4A ved stoppestedet og snakker med Henrik. Klokken er fem minutter over otte, og den næste bus går kvart over otte. Om morgenen kører bussen hvert tiende minut, og turen til banegården tager tyve minutter. Om mandagen møder Emma på sprogskolen klokken halv ni, så hun tager normalt bussen fem minutter i otte. Lige pludselig kommer bussen, og de siger god arbejdsdag til hinanden.}
\end{itemize}
\end{answerbox}
